\section{Conclusion}
\label{sec:conclusion}

The central thesis of this report is that the inspection threshold of an industrial anomaly detection system must not be treated as an arbitrary academic parameter -- it is a practical business decision that must be calibrated directly against the real consequences of false alarms and missed defects on each specific production line. A stripped fastener thread and a contaminated pharmaceutical tablet differ in their defect escape liability by orders of magnitude. Any system that applies a uniform, symmetric threshold across both settings is fundamentally miscalibrated for at least one of them. The principal achievement of our platform is making this calibration simple, transparent, and mathematically sound.

The underlying engineering architecture is purpose-built to deliver this operational vision. PatchCore utilising a frozen ResNet-18 backbone delivers a macro-average Image $F_1$-score of $0.929$ and AUPIMO of $0.603$ across all 15 MVTec AD categories without requiring complex model training on factory computers. The decoupled coreset memory bank accommodates raw material and ambient illumination drift in seconds. High-resolution spatial explainability overlays render every divert decision auditable by line-side quality inspectors. Together, these properties elevate automated visual quality assurance from an unpredictable engineering experiment into an auditable, deterministic, and highly reliable operational asset.
